% Chapter 286 reference apparatus; every relation retains its actual scope.
\chapter{Strength, Consistency, and Independence Tables}
\label{sec:p14-strength-consistency-and-independence}

These tables answer different questions: which conclusions a theory proves,
which theories it interprets, which additional set-existence principles a
mathematical theorem requires, and which extensions are relatively consistent.
They do not combine those questions into one numerical ranking. A theory
name must be read with its language, logic, axioms, schemes and presentation.
In particular, equal displayed ordinals do not identify theories, and an
implication between consistency assertions is not an implication between
their nonlogical axioms.

\begin{hbconvention}{def:p14-strength-relations}{Reading the comparison relation}
Fix the stated proof calculus for each theory. Write
\(\operatorname{Con}_{\mathrm{ext}}(T)\) for the external assertion that no
finite \(T\)-derivation concludes falsity. An internal arithmetic consistency
sentence instead uses the particular proof predicate in
Definition~\ref{def:p05-consistency-presentation}; in the selected PA
application this is the direct tag-19 predicate of
Section~\ref{sec:second-incompleteness-01}.

The following relations are kept separate.
\begin{description}
 \item[Theorem inclusion.] In a common language, every theorem of \(U\)
 is a theorem of \(V\). An axiom inclusion together with a compatible
 calculus supplies this direction. ``Strict'' additionally requires a
 sentence proved by \(V\) but not by \(U\).
 \item[Conservation.] For a specified sentence class \(\mathcal C\) of
 the old language, every \(\mathcal C\)-sentence proved by the extension
 is already proved by the base. The class restriction is part of the result.
 \item[Interpretation.] Specified translation data interpret \(U\) in
 \(V\); translated \(U\)-theorems become \(V\)-theorems. The source
 and target languages need not coincide.
 \item[Relative consistency.] A metatheorem has the direction
 \(\operatorname{Con}_{\mathrm{ext}}(V)\to
 \operatorname{Con}_{\mathrm{ext}}(U)\). A formalized version separately
 names the arithmetic base and proof presentations.
 \item[Ordinal calibration.] An analysis assigns the displayed ordinal
 under a specified notation, induction/reflection reading and formula class.
 It is not automatically an interpretation or an adjacent strict
 consistency separation.
\end{description}
These are statements about theories and proofs, not properties of one
intended model. The interpretation convention is
Definition~\ref{def:p05-relative-interpretation}; the ordinal convention is
Definition~\ref{def:p07-relative-proof-theoretic-ordinal}.
\end{hbconvention}

\section{Arithmetic and reverse-mathematical systems}
\label{sec:p14-strength-arithmetic}

The first table distinguishes first-order arithmetic theories before the
second-order set-existence schemes are compared. Scheme parameters are
universally closed; the numbered definitions, rather than an informal name,
fix their exact allowed formula classes.

\begingroup\small
\begin{longtable}{@{}p{.15\textwidth}p{.39\textwidth}p{.37\textwidth}@{}}
\caption{Arithmetic presentations and what their comparison means.}
\label{tab:p14-strength-arithmetic}\\
\toprule
Theory & Exact presentation & Comparison and statement/source destination \\
\midrule
\endfirsthead
\caption[]{Arithmetic presentations and what their comparison means (continued)}\\
\toprule
Theory & Exact presentation & Comparison and statement/source destination \\
\midrule
\endhead
\midrule\multicolumn{3}{r}{\emph{continued on next page}}\\
\endfoot
\endlastfoot
\(\mathsf Q\) & Seven number axioms; no induction scheme. &
PA extends these axioms. Use Table~\ref{tab:p05-q-pa-axioms}; essential
undecidability is a different property from the size of this axiom list,
Definition~\ref{def:p05-essential-undecidability}. \\
\addlinespace
\(\mathsf{PRA}\) & Primitive-recursive function symbols and defining
equations, with quantifier-free induction. &
The common-language comparison with \(\mathsf I\Sigma_1\) uses the
standard definitional treatment of these symbols, not raw inclusion of
different alphabets; Definition~\ref{def:p07-pra}. \\
\addlinespace
\(\mathsf I\Sigma_n\), \(n\geq1\) &
\(\mathsf Q\) plus induction restricted to \(\Sigma^0_n\) formulas. &
Every fixed fragment's induction instances are PA instances. The index
\(n\) is fixed in a fragment; the family is not one unindexed theory.
Definition~\ref{def:p07-isigma-fragments}. \\
\addlinespace
\(\PA\) & Classical first-order arithmetic; induction for every arithmetic
formula, with numerical parameters. &
The exact scheme is Table~\ref{tab:p05-q-pa-axioms}. Its standard-model
truth is not identical to its theorem set; see
Theorems~\ref{thm:godel-first} and~\ref{thm:rosser}. \\
\addlinespace
\(\mathsf{HA}\) & The stated arithmetic/equality axioms and full arithmetic
induction, with intuitionistic first-order logic. &
Definition~\ref{def:p07-heyting-arithmetic}. Its common ordinal with PA
does not turn intuitionistic proof rules into classical ones. Its
disjunction and numerical-existence metatheorems remain cited in
Theorems~\ref{thm:p07-ha-disjunction-property} and
\ref{thm:p07-ha-numerical-existence}. \\
\bottomrule
\end{longtable}
\endgroup

\begingroup\small
\begin{longtable}{@{}p{.17\textwidth}p{.37\textwidth}p{.37\textwidth}@{}}
\caption{Reverse-mathematical calibration over the fixed
\(\mathsf{RCA}_0\) base. A reversal supplies both directions.}
\label{tab:p14-strength-reverse-mathematics}\\
\toprule
System & Added principle / scheme & Representative calibration and boundary \\
\midrule
\endfirsthead
\caption[]{Reverse-mathematical calibration over the fixed
\(\mathsf{RCA}_0\) base. A reversal supplies both directions (continued)}\\
\toprule
System & Added principle / scheme & Representative calibration and boundary \\
\midrule
\endhead
\midrule\multicolumn{3}{r}{\emph{continued on next page}}\\
\endfoot
\endlastfoot
\(\mathsf{RCA}_0\) & \(\Sigma^0_1\) induction and
\(\Delta^0_1\) comprehension, with the declared set and number parameters. &
The common base, Definition~\ref{def:p05-rca0}; identifying the base is not
a reversal of a new mathematical assertion. \\
\addlinespace
\(\mathsf{WKL}_0\) & Paths through every coded infinite binary tree. &
Equivalent over the base to finite subcovers of countable rational-interval
covers of the coded \([0,1]\); Table~\ref{tab:p05-reverse-calibration},
source-cited \cite{Simpson2009}. Arbitrary uncountable covers are not the
formalized theorem in this row. \\
\addlinespace
\(\mathsf{ACA}_0\) & Arithmetical comprehension, including set parameters. &
Equivalent over the base to existence of the range of every injection
\(\N\to\N\); Table~\ref{tab:p05-reverse-calibration}, source-cited
\cite{Simpson2009}. This adds sets, not first-order quantifiers alone. \\
\addlinespace
\(\mathsf{ATR}_0\) & Arithmetical transfinite recursion along each coded
countable well-order. &
Equivalent over the base to comparison of two coded countable well-orders
by an initial-segment isomorphism; the entire target may be the segment.
Table~\ref{tab:p05-reverse-calibration}, source-cited \cite{Simpson2009}. \\
\addlinespace
\(\Pi^1_1\text{-}\mathsf{CA}_0\) & Comprehension for the displayed
universal-set-quantifier formulas. &
Equivalent over the base to forming the well-founded-tree index set for
every coded sequence of trees on \(\N^{<\N}\).
Table~\ref{tab:p05-reverse-calibration}, source-cited \cite{Simpson2009}. \\
\addlinespace
\(\mathsf Z_2\) & Full comprehension and the full two-sorted induction
scheme. &
The ambient formal theory of Definition~\ref{def:p05-z2}. Full second-order
semantics and this formal axiom system are different notions; do not
identify all its models with \((\N,\mathcal P(\N))\). \\
\bottomrule
\end{longtable}
\endgroup

The strict Big Five theorem-inclusion chain is
Theorem~\ref{thm:p05-big-five-strict}, cited to \textcite{Simpson2009}.
It does not say that the five systems exhaust the intermediate or
incomparable principles. In particular, strictness in the two-sorted
language can coexist with equality of first-order consequences.

\begingroup\small
\begin{longtable}{@{}p{.24\textwidth}p{.35\textwidth}p{.32\textwidth}@{}}
\caption{Conservation bridges, not identities of theories.}
\label{tab:p14-strength-conservation}\\
\toprule
Comparison & Exactly retained consequences & Evidence and restriction \\
\midrule
\endfirsthead
\caption[]{Conservation bridges, not identities of theories (continued)}\\
\toprule
Comparison & Exactly retained consequences & Evidence and restriction \\
\midrule
\endhead
\midrule\multicolumn{3}{r}{\emph{continued on next page}}\\
\endfoot
\endlastfoot
\(\mathsf I\Sigma_1\) over \(\mathsf{PRA}\) &
\(\Pi^0_2\)-conservation for the matched standard presentations;
the associated external consistency assertions are equivalent. &
Source-cited Parsons result, Table~\ref{tab:p06-consistency-strength-master}
and \cite{Parsons1970}. It is not unrestricted same-language conservation. \\
\addlinespace
\(\mathsf{RCA}_0\), \(\mathsf{WKL}_0\) and
\(\mathsf I\Sigma_1\) &
The first-order part of each of the two second-order systems is
\(\mathsf I\Sigma_1\). WKL adds genuinely new two-sorted consequences
without adding first-order ones. &
Source-cited comparison statements in
\cite[Remarks~I.7.6, I.10.5; outline~I.13]{Simpson2006ComparisonPreview};
full model/conservation proofs are in
\cite[Sections~IX.1--IX.3]{Simpson2009}, not in the preview. \\
\addlinespace
\(\mathsf{ACA}_0\) over \(\PA\) &
Exactly the same first-order arithmetic sentences; the standard
presentations are externally equiconsistent. &
Definition~\ref{def:p05-full-aca},
Table~\ref{tab:p06-consistency-strength-master} and
\cite[Section~IX.1]{Simpson2009}. This is not literal equality of their
languages or of their two-sorted theorem sets. \\
\bottomrule
\end{longtable}
\endgroup

The unsubscripted full-induction \(\mathsf{ACA}\) of
Definition~\ref{def:p05-full-aca} must not be substituted for
\(\mathsf{ACA}_0\): Theorem~\ref{thm:p05-ph-upper-calibration} uses that
extra induction for its Paris--Harrington upper bound. Equally, a restriction
to one first-order sentence class cannot be dropped just because the word
``conservative'' appears in a table heading.

\section{Proof-theoretic ordinals and interpretation strength}
\label{sec:p14-strength-ordinals}

The values below reproduce exact calibration claims, not a universal
definition of strength. Fix the recursive order coding, the formula class
for transfinite induction and the analysis-relative meaning of the ordinal
as in Definition~\ref{def:p07-relative-proof-theoretic-ordinal}. The ordinal
arithmetic itself is Definition~\ref{def:p06-ordinal-arithmetic}.

\begingroup\small
\begin{longtable}{@{}p{.21\textwidth}p{.18\textwidth}p{.52\textwidth}@{}}
\caption{Source-qualified ordinal calibrations and their limits.}
\label{tab:p14-strength-ordinals}\\
\toprule
Presentation & Displayed ordinal & Meaning and full-analysis destination \\
\midrule
\endfirsthead
\caption[]{Source-qualified ordinal calibrations and their limits (continued)}\\
\toprule
Presentation & Displayed ordinal & Meaning and full-analysis destination \\
\midrule
\endhead
\midrule\multicolumn{3}{r}{\emph{continued on next page}}\\
\endfoot
\endlastfoot
\(\mathsf{PRA}\) & \(\omega^\omega\) &
The modern primitive-recursive presentation and its conventional first-order
calibration, not a definition of finitism. The same displayed calibration
for \(\mathsf I\Sigma_1\) uses the Parsons comparison;
Table~\ref{tab:p07-ordinal-calibration}, \cite{Pohlers2009,Parsons1970}. \\
\addlinespace
\(\mathsf I\Sigma_n\), fixed \(n\geq1\) & \(\omega_{n+1}\) &
Here \(\omega_0=1\), \(\omega_{k+1}=\omega^{\omega_k}\).
In particular \(\omega_2=\omega^\omega\); an indexing shift changes the
printed value. Table~\ref{tab:p07-ordinal-calibration}, source-cited
\cite{Pohlers2009}. \\
\addlinespace
\(\PA\) & \(\varepsilon_0\) &
The selected Gentzen reduction/induction calibration. Exact inherited
Section~\ref{sec:normalization-curry-howard-04} is cited, with its calculus,
coding, reduction and decrease requirements exposed in
Definition~\ref{def:p07-gentzen-reduction-method};
\cite{Gentzen1936,TroelstraSchwichtenberg2000}. \\
\addlinespace
\(\mathsf{HA}\) & \(\varepsilon_0\) &
The source ordinal analysis does not identify the intuitionistic and
classical theories. Table~\ref{tab:p07-ordinal-calibration}, source-cited
\cite{Pohlers2009}. \\
\addlinespace
\(\mathsf{ACA}_0\) & \(\varepsilon_0\) &
Its recursive-well-order analysis and its first-order PA conservation
bridge are distinct supporting statements; neither says its language is
first-order. Table~\ref{tab:p07-ordinal-calibration},
\cite{Simpson2009,Pohlers2009}. \\
\addlinespace
\(\mathsf{ATR}_0\) & \(\Gamma_0\) &
The second-order calibration. The separate autonomous predicative
progressions also display \(\Gamma_0\), but that shared ordinal is not a
literal identity of progressions and ATR.
Theorem~\ref{thm:p07-feferman-schuette} and
Table~\ref{tab:p07-ordinal-calibration},
\cite{Feferman1964,Schuette1977,Simpson2009}. \\
\bottomrule
\end{longtable}
\endgroup

An ordinal bound for a particular proof procedure may be deliberately
nonminimal. For example, the fully proved finite-set propositional Cut
strategy uses a decreasing measure below \(\omega^{\omega^2}\),
Theorem~\ref{thm:p07-prop-cut-termination} and
Annex~\ref{annex:p7.arithmetic.1}. That is a measure on a specified
derivation grammar and strategy, not an asserted proof-theoretic ordinal
for propositional logic and not Gentzen's arithmetic assignment.

\begingroup\small
\begin{longtable}{@{}p{.24\textwidth}p{.34\textwidth}p{.33\textwidth}@{}}
\caption{Interpretation and extraction comparisons.}
\label{tab:p14-strength-interpretations}\\
\toprule
Input / target & What is transferred & Conditions and evidence \\
\midrule
\endfirsthead
\caption[]{Interpretation and extraction comparisons (continued)}\\
\toprule
Input / target & What is transferred & Conditions and evidence \\
\midrule
\endhead
\midrule\multicolumn{3}{r}{\emph{continued on next page}}\\
\endfoot
\endlastfoot
An interpretation of \(U\) in \(V\) &
\(U\derives\sigma\) implies \(V\derives\sigma^{\mathcal I}\);
\(\operatorname{Con}_{\mathrm{ext}}(V)\) implies
\(\operatorname{Con}_{\mathrm{ext}}(U)\). &
Nonempty definable domain, equality/congruence, function-graph totality and
uniqueness modulo equality, and translated axioms;
Theorem~\ref{thm:p05-interpretation-transfer}, cited \cite{TMR1953}. \\
\addlinespace
\(\mathsf Q\) in \(\PA\) & Identity interpretation on the arithmetic
language. & Example~\ref{def:p05-interpretation-examples} supplies the
exact elementary example; no reverse interpretation or unrestricted
conservation follows from that example. \\
\addlinespace
A definitional expansion and its base &
Replace the new vocabulary by its defining formulas to interpret the
expansion back into its base. &
The actual defining axioms and unique-value conditions are required;
Example~\ref{def:p05-interpretation-examples}. This differs from an
arbitrary extension by new existence axioms. \\
\addlinespace
HA through Dialectica &
A closed HA proof yields closed System-T witness terms and a proof of the
translated matrix in \(\mathsf T_{\mathrm{qf}}\). &
The exact finite-type target and translation are
Definitions~\ref{def:p07-dialectica-target} and
\ref{def:p07-dialectica-translation}; soundness and the relative consistency
consequence are cited in Theorems~\ref{thm:p07-dialectica-soundness} and
\ref{thm:p07-dialectica-consistency}, \cite{AvigadFeferman1998}.
Term extraction is not formula equivalence in HA. \\
\bottomrule
\end{longtable}
\endgroup

For a quotient interpretation, definable unique representatives are not
assumed: uniqueness is modulo the interpreted equality relation. Nor does
an effective interpretation transfer plain undecidability without further
hypotheses. The cited essential-undecidability transfer additionally requires
a consistent c.e. target and an essentially undecidable source.

\section{Set-theoretic axioms and large-cardinal comparisons}
\label{sec:p14-strength-set-theory}

All class-theory names use the choices declared in
Table~\ref{tab:p06-set-theory-variants}. Here modern \(\mathsf Z\) retains
Foundation but omits Replacement. The handbook's \(\mathsf{NBG}\) includes
Global Choice; its \(\mathsf{GB}\) designation does not. Global Choice for
\(\mathsf{MK}\) is written as an explicit addition when a comparison needs
it. These choices must not be silently changed to another author's naming
convention.

\begingroup\small
\begin{longtable}{@{}p{.25\textwidth}p{.34\textwidth}p{.32\textwidth}@{}}
\caption{Set-theory comparisons by relation, not by a single ladder.}
\label{tab:p14-strength-set-axioms}\\
\toprule
Compared theories / structures & Supported relation & Restriction and evidence \\
\midrule
\endfirsthead
\caption[]{Set-theory comparisons by relation, not by a single ladder (continued)}\\
\toprule
Compared theories / structures & Supported relation & Restriction and evidence \\
\midrule
\endhead
\midrule\multicolumn{3}{r}{\emph{continued on next page}}\\
\endfoot
\endlastfoot
\(\mathsf Z\), \(\mathsf{ZF}\) & Adding Replacement gives the declared
ZF presentation. & Literal axiom extension in this convention; no adjacent
strict consistency theorem is inferred solely from the added scheme.
Table~\ref{tab:p06-set-theory-variants}, \cite{Jech2003}. \\
\addlinespace
\(\mathsf{ZF}\), \(\mathsf{ZFC}\) & Adding AC gives ZFC;
constructibility supplies external equiconsistency. &
ZF does not thereby prove Choice. The positive constructibility implication
is Theorem~\ref{thm:ch-relative-consistency}; its proof uses the cited
\(L\)-axiom-relativization scheme, \cite{Godel1938,Jech2003}. \\
\addlinespace
\(\mathsf{NBG}\), \(\mathsf{ZFC}\) &
The declared Global-Choice NBG is conservative for set-language sentences. &
Class quantifiers are not retained in the conserved sentence class;
Table~\ref{tab:p06-set-theory-variants}, source-cited
\cite{Mendelson1997,Jech2003}. \\
\addlinespace
\(\mathsf{GB}\), \(\mathsf{MK}\) & Predicative versus impredicative
class comprehension; the latter allows class quantifiers in its scheme. &
Compare matched set axioms and choice additions before asserting a
conservation or consistency-strength edge. The class-theory distinction is
Table~\ref{tab:p06-set-theory-variants}, \cite{Kelley1975}; no identity with
ZFC plus one inaccessible is asserted. \\
\addlinespace
\(\mathsf{KP}\), \(\mathsf{ZF}\) &
KP restricts Separation/Collection to bounded formulas and has no Power Set
axiom; the exact KP-with-Power-Set variant is separate. &
The axiom-profile comparison is Table~\ref{tab:p06-set-theory-variants},
\cite{Barwise1975}. A guessed common ordinal or model-existence equivalence
is not supplied by these differences. \\
\addlinespace
\((\omega,\mathcal P(\omega))\) in \(\mathsf Z\) &
The cited standard-model route compares full second-order arithmetic with
set theory. & Table~\ref{tab:p06-consistency-strength-master},
\cite{Jech2003,Simpson2009}. Its set-existence/model route is not a
literal inclusion of two different languages. \\
\bottomrule
\end{longtable}
\endgroup

\begingroup\small
\begin{longtable}{@{}p{.19\textwidth}p{.33\textwidth}p{.39\textwidth}@{}}
\caption{Conditional large-cardinal implications. Each implication runs
from the stated stronger property to the stated weaker one.}
\label{tab:p14-strength-large-cardinals}\\
\toprule
Stronger hypothesis & Conclusion at the same cardinal & Source and what is not inferred \\
\midrule
\endfirsthead
\caption[]{Conditional large-cardinal implications. Each implication runs
from the stated stronger property to the stated weaker one (continued)}\\
\toprule
Stronger hypothesis & Conclusion at the same cardinal & Source and what is not inferred \\
\midrule
\endhead
\midrule\multicolumn{3}{r}{\emph{continued on next page}}\\
\endfoot
\endlastfoot
\(\kappa\) supercompact & \(\kappa\) measurable. &
Closure of the embedding target is part of supercompactness,
Definition~\ref{def:p06-supercompact-cardinal}; implication is cited in
Theorem~\ref{thm:p06-large-cardinal-implications}, \cite{Kanamori2003}. \\
\addlinespace
\(\kappa\) measurable & \(\kappa\) weakly compact. &
The witness is a nonprincipal \(\kappa\)-complete ultrafilter, not ordinary
real-measure measurability; Definition~\ref{def:p06-measurable-cardinal},
Theorem~\ref{thm:p06-large-cardinal-implications}, cited \cite{Kanamori2003}. \\
\addlinespace
\(\kappa\) weakly compact & \(\kappa\) Mahlo. &
The declared partition definition is
Definition~\ref{def:p06-weakly-compact-cardinal}. Its tree formulation also
requires strong inaccessibility, Theorem~\ref{thm:p06-weak-compact-tree}.
Both implications are source-cited \cite{Kanamori2003}. \\
\addlinespace
\(\kappa\) Mahlo & \(\kappa\) strongly inaccessible. &
Use the strong Mahlo convention of Definition~\ref{def:p06-mahlo-cardinal},
not an unannounced weak-Mahlo variant;
Theorem~\ref{thm:p06-large-cardinal-implications}, cited \cite{Kanamori2003}. \\
\addlinespace
\(\kappa\) strongly inaccessible &
\((V_\kappa,\in)\) is a set model of ZFC. &
Theorem~\ref{thm:p06-inaccessible-stage}, cited \cite{Kanamori2003}.
The existence of this cardinal is an extra hypothesis, not a ZFC existence
theorem or an assumption derived from mere external consistency. \\
\bottomrule
\end{longtable}
\endgroup

The large-cardinal theorem gives one-way comparisons between existence
extensions of ZFC. It does not by itself prove strictness of every adjacent
consistency strength. Measurability excludes \(V=L\),
Theorem~\ref{thm:p06-scott-measurable}; this is not a prohibition on all
large cardinals in \(L\). Inaccessible, Mahlo and weakly compact are the
separately cited compatible rungs in
Table~\ref{tab:p06-large-cardinal-ladder}.

An inaccessible stage is sufficient for the indicated set model and the
external countable-transitive-model construction; it is not equivalent to an
arbitrary transitive model. Keep this distinction from
Warning~\ref{def:p06-ctm-inaccessible-boundary} when reading the forcing
examples. Cardinal arithmetic inside a ground model, transitivity, external
countability and an internally computed chain condition must each retain
their superscripts and quantifier domains.

\section{Choice principles, relative consistency, and independence}
\label{sec:p14-strength-choice-independence}

\begingroup\small
\begin{longtable}{@{}p{.26\textwidth}p{.34\textwidth}p{.31\textwidth}@{}}
\caption{Choice comparisons over ZF. Different compactness formulations
consume different principles.}
\label{tab:p14-strength-choice}\\
\toprule
Principle / implication & Exact statement or restriction & Evidence and converse boundary \\
\midrule
\endfirsthead
\caption[]{Choice comparisons over ZF. Different compactness formulations
consume different principles (continued)}\\
\toprule
Principle / implication & Exact statement or restriction & Evidence and converse boundary \\
\midrule
\endhead
\midrule\multicolumn{3}{r}{\emph{continued on next page}}\\
\endfoot
\endlastfoot
AC, well-ordering, Zorn, Hausdorff maximality, Tukey, arbitrary nonempty
products & Equivalent full-Choice forms over ZF; Zorn's upper-bound condition
includes the empty chain. &
The full equivalence list is cited in Theorem~\ref{thm:p06-zorn}; the complete
AC-to-Zorn proof is Theorem~\ref{thm:p06-ac-to-zorn} and
Annex~\ref{annex:p6.choice.1}. \\
\addlinespace
\(\mathsf{AC}\to\mathsf{DC}\to\mathsf{AC}_\omega
\to\mathsf{AC}_\omega^{\mathrm{fin}}\) &
Dependent Choice starts at a specified point of a serial relation; Countable
Choice selects from a given sequence of nonempty sets. &
Definitions~\ref{def:p06-countable-dependent-choice} and cited
Theorem~\ref{thm:p06-choice-implications}, \cite{HowardRubin1998}.
The reverse implications fail in the stated source model comparisons. \\
\addlinespace
Rasiowa--Sikorski with a prescribed start & An arbitrary nonempty poset
and a countable family of dense sets; successive extensions depend on
the previous condition. & The actual \(\mathsf{ZF}+\mathsf{DC}\)
construction is Lemma~\ref{thm:p10-rasiowa-sikorski} and
Annex~\ref{annex:p10.rs.1}. A fixed enumeration of the poset permits
least-index choices in \(\mathsf{ZF}\); a finite dense family needs
no DC. The displayed sufficient use is not a claim of equivalence to DC. \\
\addlinespace
\(\mathsf{AC}\to\mathsf{BPI}\to
\mathsf{AC}_\omega^{\mathrm{fin}}\) & BPI is Boolean prime ideal /
ultrafilter extension, Definition~\ref{def:ultrafilter-lemma}. &
Theorem~\ref{thm:p06-choice-implications}, cited
\cite{HowardRubin1998,HalpernLevy1971}. BPI does not imply full AC; no
DC--BPI comparison follows from the displayed chains. \\
\addlinespace
Recent BPI proof and compactness development & \(\mathsf{BPI}\),
ultrafilter extension and unrestricted propositional compactness,
with filters proper and Boolean algebras set-sized. &
\cite{Repicky2015BPI} is the separately identified 2015 source.
The supplied handbook constructions are
Theorems~\ref{thm:p01-bpi-equivalence}
and~\ref{thm:p01-bpi-propositional-compactness};
the ND/free-Boolean construction is
Theorem~\ref{thm:p02-fbc-completeness}. Neither proves BPI in ZF
or the Halpern--Levy separation construction. \\
\addlinespace
Arbitrary propositional compactness & Variables may form an arbitrary set;
equivalent over ZF to BPI. &
The full matched result is Theorem~\ref{thm:p02-arbitrary-prop-compactness-bpi},
not the countable binary-tree argument. \\
\addlinespace
Countable binary-tree compactness & A subtree of \(2^{<\omega}\) with a
node at every level has a branch, in classical ZF without Choice. &
Theorem~\ref{thm:p06-konig-infinity} and complete
Annex~\ref{annex:p6.choice.2}; arbitrary unordered finite successor sets
are not this fixed binary presentation. \\
\addlinespace
Tychonoff variants & Arbitrary compact spaces give the full-Choice
equivalence; compact Hausdorff spaces give BPI instead. &
Theorem~\ref{thm:p06-zorn} and Table~\ref{tab:p06-ac-equivalents},
source-cited \cite{Jech1973Choice,HowardRubin1998}. The Hausdorff qualification
cannot be dropped from the weaker row. \\
\bottomrule
\end{longtable}
\endgroup

The choice-free binary-tree theorem uses classical ZF as its base.
\(\mathsf{WKL}_0\) instead adds the coded arithmetical version over the
much weaker \(\mathsf{RCA}_0\). The different bases explain why a ZF
proof without Choice does not make weak K\"onig's lemma a theorem of
\(\mathsf{RCA}_0\).

\begin{hbtheorem}{thm:p14-strength-proof-transfer}{proved=annex:p14.strength.proof-transfer}{What a contradiction-preserving proof transfer establishes}
Let \(U,V\) have specified calculi, and suppose a transformation sends
every finite \(U\)-proof of falsity to a finite \(V\)-proof of falsity.
Then
\[
 \operatorname{Con}_{\mathrm{ext}}(V)\longrightarrow
 \operatorname{Con}_{\mathrm{ext}}(U).
\]
For the formalized assertion, let \(\mathsf B\) be a theory with first-order
arithmetic language and logic. Fix proof formulas \(P_U(p,x),P_V(q,x)\),
the respective numerical falsity codes \(c_U,c_V\), and a transformation
graph formula \(F(p,q)\). If \(\mathsf B\) proves
\[
 \forall p\,\exists q\,F(p,q),\qquad
 \forall p\forall q\,
 ((P_U(p,\bar c_U)\land F(p,q))\to P_V(q,\bar c_V)),
\]
then it proves
\[
 \neg\exists q\,P_V(q,\bar c_V)\ \to\
 \neg\exists p\,P_U(p,\bar c_U).
\]
Here the displayed bars are numerals for the fixed codes. External existence
of a transformation alone does not assert that \(\mathsf B\) proves these
totality and correctness formulas.
\end{hbtheorem}

\begin{hbtheorem}{thm:p14-strength-independence-transfer}{proved=annex:p14.strength.independence-transfer}{The two-extension route to independence}
Let \(T\) use classical logic, let \(\varphi\) be a sentence, and assume
both metatheoretic implications
\[
\begin{aligned}
 \operatorname{Con}_{\mathrm{ext}}(T)&\to
       \operatorname{Con}_{\mathrm{ext}}(T+\varphi),\\
 \operatorname{Con}_{\mathrm{ext}}(T)&\to
       \operatorname{Con}_{\mathrm{ext}}(T+\neg\varphi).
\end{aligned}
\]
If \(T\) is consistent, it proves neither \(\varphi\) nor
\(\neg\varphi\). The implication for \(T+\neg\varphi\) excludes
\(T\derives\varphi\); the implication for \(T+\varphi\) excludes
\(T\derives\neg\varphi\). A single direction is not a two-sided
independence theorem.
\end{hbtheorem}

\begingroup\small
\begin{longtable}{@{}p{.14\textwidth}p{.34\textwidth}p{.43\textwidth}@{}}
\caption{Set-theoretic independence: both relative-consistency directions
and their actual support. Every consistency assertion in this table is external.}
\label{tab:p14-strength-set-independence}\\
\toprule
Statement / base & Positive and negative extensions & Method, assumptions and statement/proof links \\
\midrule
\endfirsthead
\caption[]{Set-theoretic independence: both relative-consistency directions
and their actual support. Every consistency assertion in this table is external (continued)}\\
\toprule
Statement / base & Positive and negative extensions & Method, assumptions and statement/proof links \\
\midrule
\endhead
\midrule\multicolumn{3}{r}{\emph{continued on next page}}\\
\endfoot
\endlastfoot
CH over ZFC & Consistency of \(\mathsf{ZFC}\) implies consistency of
\(\mathsf{ZFC}+\mathsf{CH}\) and of
\(\mathsf{ZFC}+\neg\mathsf{CH}\). & Positive:
Theorem~\ref{thm:ch-relative-consistency} has the
relativization derivation, using the cited constructibility scheme.
Negative: the cited Boolean-valued theorem in
Section~\ref{sec:constructibility-forcing-02}, \cite{Cohen1963,Jech2003}.
The complete final nonderivability argument is
Theorem~\ref{thm:ch-nonderivability}; it does not reprove the forcing import. \\
\addlinespace
GCH over ZFC & Consistency of \(\mathsf{ZFC}\) implies consistency of
\(\mathsf{ZFC}+\mathsf{GCH}\) and of
\(\mathsf{ZFC}+\neg\mathsf{GCH}\). & The \(L\) direction is included in
Theorem~\ref{thm:ch-relative-consistency}. The negative direction follows
from the cited not-CH extension since GCH implies CH; the wider regular-
cardinal pattern discussion is Table~\ref{tab:p06-independence-catalog},
\cite{Easton1970}. No arbitrary singular-cardinal freedom is asserted. \\
\addlinespace
AC over ZF & Consistency of \(\mathsf{ZF}\) implies consistency of
\(\mathsf{ZF}+\mathsf{AC}\) and of
\(\mathsf{ZF}+\neg\mathsf{AC}\). & Positive: the constructible
inner-model implication.
Negative: symmetric/permutation methods with the cited pure-set transfer,
Table~\ref{tab:p06-independence-catalog}, \cite{Jech1973Choice,Halbeisen2017}.
Theorem~\ref{thm:p06-basic-fraenkel-model} is a ZFA atom-model example,
not already a pure ZF model. \\
\bottomrule
\end{longtable}
\endgroup

The external countable transitive ground model used in a pedagogical forcing
construction is a stronger input than mere consistency. A proof that starts
with that model has not thereby established an implication from
\(\operatorname{Con}_{\mathrm{ext}}(\mathsf{ZFC})\) alone. The exact
distinction is stated in Section~\ref{sec:constructibility-forcing-02};
the Boolean-valued negative-CH import expressly removes that extra model
hypothesis. Likewise the fully proved continuum calculation
Theorem~\ref{thm:p06-many-cohen-continuum},
Annex~\ref{annex:p6.force.11}, keeps its ground-model equation
\(\kappa^{\aleph_0}=\kappa\) and the cited forcing theorem as prerequisites.

\begingroup\small
\begin{longtable}{@{}p{.18\textwidth}p{.26\textwidth}p{.47\textwidth}@{}}
\caption{Arithmetic nonderivability versus two-sided independence.}
\label{tab:p14-strength-arithmetic-independence}\\
\toprule
Sentence / base & Required metatheory & Exact conclusion and evidence \\
\midrule
\endfirsthead
\caption[]{Arithmetic nonderivability versus two-sided independence (continued)}\\
\toprule
Sentence / base & Required metatheory & Exact conclusion and evidence \\
\midrule
\endhead
\midrule\multicolumn{3}{r}{\emph{continued on next page}}\\
\endfoot
\endlastfoot
\(G_T\), c.e. \(T\supseteq\mathsf Q\) & Consistency for the positive
sentence; \(\Sigma^0_1\)-soundness or omega-consistency for its negation. &
\(T\notderives G_T\); the additional hypothesis also gives
\(T\notderives\neg G_T\). Theorem~\ref{thm:godel-first} explicitly has
proof-sketch standing; its fixed-point complexity/coding prerequisites
are not supplied by this table. \\
\addlinespace
\(R_T\), c.e. \(T\supseteq\mathsf Q\) & Consistency alone. &
Both signs are unprovable. Theorem~\ref{thm:rosser} explicitly has
proof-sketch standing with the bounded-comparison steps specified;
\cite{Rosser1936}. \\
\addlinespace
Selected \(\mathsf{Con}(T)\) & Consistent c.e. arithmetic extension of
\(\mathsf Q\); its chosen predicate satisfies D1--D3. &
\(T\notderives\mathsf{Con}(T)\), not automatically two-sided
independence. Theorem~\ref{thm:godel-second} gives the complete conditional
argument. The selected PA D1--D3 proofs are
Theorems~\ref{thm:p05-direct-pa-internal-mp} and
\ref{thm:p05-direct-pa-d3}. \\
\addlinespace
Paris--Harrington, PA & Consistency for the cited unprovability result. &
True in the standard model and unprovable in PA;
Theorem~\ref{thm:p05-paris-harrington}, cited
\cite{ParisHarrington1977}. Its full-ACA upper bound is
Theorem~\ref{thm:p05-ph-upper-calibration}; bare ACA-zero does not suffice. \\
\addlinespace
Goodstein universal termination, PA & Consistency for the cited
unprovability result; the stated hereditary-base schedule. &
True universal termination and PA-nonderivability are separate
Theorems~\ref{thm:p05-goodstein-termination} and
\ref{thm:p05-goodstein-independence}, cited
\cite{Goodstein1944,KirbyParis1982}. No generic changed-base variant is
silently substituted. \\
\addlinespace
Kirby--Paris hydra, PA & Exact stage-counted regrowth; recursive strategies
in the first-order unprovability formulation. &
The external all-battles termination statement and the unprovable
recursive-strategy assertion have different quantifier domains;
Theorem~\ref{thm:p05-hydra}, cited \cite{KirbyParis1982}. \\
\bottomrule
\end{longtable}
\endgroup

\section{Conditions on every comparison and links to full proofs}
\label{sec:p14-strength-conditions}
\subsection{Consistency and proof-theoretic strength}
\label{sec:p14-strength-table}

There is no single unconditional total ordering containing every foundational
theory. The following chains use the indicated interpretation or proof-theoretic
comparison; arrows are not set inclusion unless the row says so.

Table~\ref{tab:p06-consistency-strength-master} and
Figure~\ref{fig:p06-consistency-strength} are the master strength surfaces.
They keep the arithmetic and reverse-mathematics tracks separate, state the
PA/ACA-zero conservation bridge, distinguish set-theory extension from
equiconsistency, and label every large-cardinal edge conditionally.  An ordinal
analysis, an interpretation, a conservation theorem, and a relative
consistency theorem are four different comparison relations; neither master
surface converts one into another without an explicit theorem.


To apply a cell, first identify its relation, then check its prerequisites
and follow its statement/proof destination. The word ``cited'' means that
the named source carries the indicated proof. A full proof of a downstream
transfer does not convert a cited premise into a locally supplied proof.

\begingroup\small
\begin{longtable}{@{}p{.24\textwidth}p{.35\textwidth}p{.32\textwidth}@{}}
\caption{Comparison conditions and diagnostic failures.}
\label{tab:p14-strength-conditions}\\
\toprule
Check & Required data / usable conclusion & Failure if omitted \\
\midrule
\endfirsthead
\caption[]{Comparison conditions and diagnostic failures (continued)}\\
\toprule
Check & Required data / usable conclusion & Failure if omitted \\
\midrule
\endhead
\midrule\multicolumn{3}{r}{\emph{continued on next page}}\\
\endfoot
\endlastfoot
Theory and logic & Same language for literal theorem inclusion, or an explicit
translation between different languages/logics. & A classical theorem need
not be a theorem of the intuitionistic system with matching arithmetic axioms. \\
\addlinespace
Scheme and parameters & Allowed formula class, number/set parameters,
universal closure and any full-induction addition. & ACA and ACA-zero,
or a fixed induction fragment and all PA induction, become wrongly identified. \\
\addlinespace
Conservation class & The exact old-language sentence class in the theorem. &
Parsons \(\Pi^0_2\) conservation is promoted to unrestricted conservation. \\
\addlinespace
Internal consistency & Axiom enumeration, proof predicate, falsity code,
numbering and formal base; D1--D3 where used. & An extensional standard-model
meaning is mistaken for provable equivalence of internal formulas;
Definition~\ref{def:p05-consistency-presentation}. \\
\addlinespace
Interpretation domain & Proven nonemptiness, equality/congruence, graph
totality and unique values modulo the interpreted equality. & A translated
axiom is evaluated on an empty domain or ill-defined quotient functions. \\
\addlinespace
Ordinal comparison & Notation algorithms, formula class, reduction notion,
decrease cases and required induction/reflection. & A schematic decreasing
number or an ordinal name is presented as a complete consistency analysis;
Definition~\ref{def:p07-gentzen-reduction-method}. \\
\addlinespace
Model-existence input & Ordinary, omega-, transitive or rank-initial models
as actually required; external countability where used. & A countable
transitive ground-model assumption is weakened to bare consistency without
an additional argument. \\
\addlinespace
Cardinal conventions & Strong inaccessible/Mahlo definition, regularity,
closure, and ground-model versus extension computation. & The tree property
alone becomes weak compactness, or conditional existence becomes a ZFC theorem. \\
\addlinespace
Independence directions & Both extension consistency results, followed by
the appropriate contradiction argument. & One-sided unprovability is labelled
independence; use Theorem~\ref{thm:p14-strength-independence-transfer}. \\
\addlinespace
Proof destination & An exact result and its matched proof or its actual
cited source, with the same hypotheses. & A different calculus, source
version, quantified claim or proof sketch is silently substituted. \\
\bottomrule
\end{longtable}
\endgroup

The two elementary transfer proofs are
Annexes~\ref{annex:p14.strength.proof-transfer} and
\ref{annex:p14.strength.independence-transfer}. For complete mathematical
constructions available at the linked destinations, use the selected PA
proof-predicate conditions, AC-to-Zorn and binary-tree proofs, the final CH
nonderivability argument, and the stated ccc/name-counting consequences.
For the deeper conservation, reversal, ordinal-analysis, inner-model and
large-cardinal comparisons, the tables retain their explicit source-relative
proof standing. This separation allows a result to be used with its actual
premises without confusing a concise lookup with the full argument.
